\section{Longest Palindrome}
\label{sec:cses-longest-palindrome}

\problemstatement{Given a string, find its longest substring that is a
palindrome (reads the same forwards and backwards). Output that
substring.}

\begin{patternbox}
\emph{``Longest palindromic substring''} is exactly what \textbf{Manacher's
algorithm} computes in $O(n)$: for every possible centre it finds the
radius of the longest palindrome centred there, reusing already-known
radii via a maintained ``current rightmost palindrome.''
\end{patternbox}

\subsection*{Manacher in one idea}

For each centre, a palindrome extends symmetrically until the characters
disagree. The naive approach re-expands every centre from scratch
($O(n^2)$). Manacher keeps the palindrome with the farthest right edge
seen so far, centred at $c$ with right boundary $r$. For a new centre $i<r$,
its mirror $2c-i$ already has a known radius, which can seed $i$'s radius
(bounded by $r-i$) before any character comparison -- so each character is
matched only a constant number of times amortised. Handling even-length
palindromes is done either with a second pass or by inserting separators.

\begin{ideabox}[Mirror Symmetry Seeds Each New Centre]
Inside the current rightmost palindrome, position $i$ mirrors $2c-i$, so
its radius starts at least as large as the mirror's (capped by the
boundary). Only growth beyond the boundary requires real comparisons,
which advance $r$ -- giving linear total work.
\end{ideabox}

\subsection*{Algorithm}

\begin{enumerate}
  \item Transform $s$ to $t$ by inserting a sentinel (e.g.\ \texttt{\#})
        between every character and at the ends, so all palindromes become
        odd-length.
  \item Compute radii $P[i]$ over $t$ using the mirror-seed rule and a
        maintained $(c,r)$.
  \item The maximum $P[i]$ locates the longest palindrome; map back to
        indices in $s$ and output the substring.
\end{enumerate}

\subsection*{Dry run}

\dryrun{$s=$``aybabtu''.}

\begin{itemize}
  \item Centred at ``bab'', the palindrome extends to radius covering
        ``bab''.
  \item No longer palindrome exists, so the answer is ``bab'' (length
        $3$).
\end{itemize}

\subsection*{C++ implementation}

\begin{lstlisting}
#include <bits/stdc++.h>
using namespace std;

int main() {
    string s;
    cin >> s;
    // transform: ^ # a # b # ... $   (sentinels avoid bounds checks)
    string t = "^";
    for (char c : s) { t += '#'; t += c; }
    t += "#$";

    int n = t.size();
    vector<int> P(n, 0);
    int c = 0, r = 0, bestLen = 0, bestCenter = 0;
    for (int i = 1; i < n - 1; i++) {
        if (i < r) P[i] = min(r - i, P[2 * c - i]);   // mirror seed
        while (t[i + P[i] + 1] == t[i - P[i] - 1]) P[i]++;
        if (i + P[i] > r) { c = i; r = i + P[i]; }
        if (P[i] > bestLen) { bestLen = P[i]; bestCenter = i; }
    }
    int start = (bestCenter - bestLen) / 2;         // map back to s
    cout << s.substr(start, bestLen) << "\n";
    return 0;
}
\end{lstlisting}

\begin{complexitybox}
\textbf{Time Complexity.} $O(n)$ -- each character is matched a constant
number of times amortised. \textbf{Space Complexity.} $O(n)$ for the
transformed string and radius array.
\end{complexitybox}

\begin{takeawaybox}
Manacher computes every palindromic radius in linear time by seeding new
centres from their mirrors within the current rightmost palindrome. The
sentinel transform unifies odd and even cases. It is the definitive tool
for palindrome-substring problems, reused in the next section to count
them all.
\end{takeawaybox}
